\section{Abstract}\label{abstract}

How long a technological civilization remains active, and what determines whether it collapses or persists, is central to assessing the prevalence of observable technosignatures in the galaxy and to projecting humanity's own long-term trajectory. We model collapse-recovery dynamics across ten plausible trajectories for Earth-originating civilization using an exploratory hybrid deterministic-stochastic simulation over a 1000-year window. We define the duty cycle of a civilization as the fraction of its total lifespan during which it remains technologically active, and find values ranging from $\sim$0.38 to 1.00, with outcomes shaped by the interplay of governance structure, resource pressure, and hazard exposure. From the duty cycle we derive an effective detectability duration that accounts for intermittent civilizational activity, and argue that intermittency systematically reduces the number of detectable civilizations relative to models that assume continuous emission. Within the assumptions of our Earth-derived scenario set, low duty cycles thus offer one possible contributing explanation, alongside others, for the apparent absence of detected technosignatures. Sensitivity analysis, including a global variance-based assessment, identifies the rate of resource depletion and the post-collapse recovery fraction as the parameters with the greatest influence on civilizational persistence across most scenarios. We discuss implications for technosignature search strategies and for the assessment of Earth's civilizational resilience.

\vspace{1em}
\noindent\textbf{Keywords:} technosignatures; duty cycle; technosphere; civilizational longevity; astrobiology
